\section{Which Models Were Evaluated?}
\label{app:model-suite-selection}
\label{sec:app-model-suite-selection}

This appendix explains the model suite used to evaluate Sparse Readout Prism:
the Qwen model suite \citep{qwen3technical,qwen35collection}, the Gemma model
suite \citep{gemma4official,gemma4collection}, Ministral
\citep{liu2026ministral3}, and targeted DeepSeek-R1 distilled comparison rows
\citep{deepseekai2025deepseekr1}.
Appendices~\ref{app:readout-factorizer-sweeps}--%
\ref{app:readout-factorizer-diagnostics-reproducibility} give the metric
summary: they report the trained widths, native $k$ settings, rowEV, top1, KL,
and usage statistics for each readout SAE setting. Here we motivate the model
choices.

Sparse Readout Prism factorizes the matrix used to score output tokens, so the
model suite is organized around readout fidelity, scale, and portability. The
chosen suite lets us compare model size and readout geometry while preserving
enough budget for tokenizer, row-norm, special-token, and matched-control
analyses.

\subsection{Selection Rationale}
\label{sec:app-model-suite-criteria}

The Qwen3.5 rows provide a fixed qualitative display setting and Qwen-family
companion analyses. Qwen3.5-9B is the capacity anchor: it is large enough to be a
meaningful readout target while still supporting recipe and width sweeps.
Qwen3.5-0.8B and Qwen3.5-2B test the same recipe at smaller unembedding
dimensions and cheaper rerun settings.

The Gemma rows provide a second tied-readout family with a different tokenizer
and final-logit softcap. They show the readout SAE recipe under
a distinct readout geometry; the softcap-correct readout evaluation is reported
in Appendix~\ref{app:readout-factorizer-transfer-comparisons}.

Additional targeted rows broaden the readout-geometry coverage.
Ministral-3-8B-Base tests transfer to an untied, softcap-free 8B-scale readout
outside Qwen. R1-Distill-Qwen-7B and R1-Distill-Llama-8B evaluate the same
recipe-selection measurements on reasoning-distilled comparison rows, making row
reconstruction and LM-head replacement fidelity visible as separate result axes.

\subsection{Final Suite}
\label{sec:app-model-suite-chosen}

\begin{table*}[!tbp]
\caption{
Evaluated model suite and role. Appendices~\ref{app:readout-factorizer-sweeps}--%
\ref{app:readout-factorizer-diagnostics-reproducibility} report the exact
dictionary settings and metrics.
}
\label{tab:app-model-suite-current}
\centering
\footnotesize
\setlength{\tabcolsep}{4pt}
\renewcommand{\arraystretch}{1.10}
\begin{tabularx}{0.94\textwidth}{@{}p{0.20\textwidth}Yp{0.22\textwidth}@{}}
\toprule
Model(s) & Why included & Role in study \\
\midrule
\shortstack[l]{\texttt{Qwen3.5-0.8B}\\\texttt{Qwen3.5-2B}} &
Small-model transfer and lower-cost repeats under the native $k=128$ and
$k=256$ regimes. &
Qwen transfer evidence. \\
\addlinespace
\texttt{Qwen3.5-9B} &
Scale anchor for recipe selection, capacity sweeps, and high-fidelity
readout replacement. &
Qwen fidelity evidence. \\
\addlinespace
\shortstack[l]{\texttt{Gemma-4-E2B}\\\texttt{Gemma-4-E4B}} &
Cross-family evidence for whether the readout SAE recipe transfers beyond
Qwen. &
Cross-family transfer evidence. \\
\addlinespace
\shortstack[l]{\texttt{R1-Distill-}\\\texttt{Qwen-7B}} &
Qwen-derived post-training comparison with an untied, softcap-free readout. &
rowEV vs.\ fidelity separation evidence. \\
\addlinespace
\shortstack[l]{\texttt{R1-Distill-}\\\texttt{Llama-8B}} &
Reasoning-distilled comparison from a non-Qwen base family (both dictionary
settings) for cross-architecture confirmation of the rowEV / fidelity gap. &
Cross-architecture post-training evidence. \\
\addlinespace
\shortstack[l]{\texttt{Ministral-3-}\\\texttt{8B-Base}} &
Cross-family 8B-scale comparison with an untied, softcap-free readout
(adds an 8B-scale non-Qwen readout geometry). &
Cross-family fidelity comparison. \\
\bottomrule
\end{tabularx}
\end{table*}

The qualitative feature analyses use one fixed low-error display setting so that
examples share the same readout SAE and reporting convention. This is an
evidence-driven display choice: the additional rows show how the same metrics
behave under different readout geometries and post-training regimes.

\subsection{Interpretation Details}
\label{sec:app-model-suite-boundaries}

The suite is an evidence-driven design for readout-factorization evaluation.
We report token frequency, token length, special-token behavior, row norm, tied
embedding status, softcap handling, and matched-target identity bias so
single-token logit decompositions can be compared on the same footing.

The evaluation makes three commitments. First, model-specific readout layers
are handled before reporting decision fidelity. Second, the SAE architecture,
training recipe, seed convention, dictionary settings, and evaluation settings
are reported with the sweep records in
Appendices~\ref{app:readout-factorizer-sweeps}--%
\ref{app:readout-factorizer-diagnostics-reproducibility}.
Third, primary reported claims are stated at the scalar-readout-target level,
especially token logits and contrastive logit differences.
